% ***************************************************
% Abstract
% ***************************************************
% TO PRODUCE A STAND-ALONE PDF OF YOUR ABSTRACT, uncomment this section and the \end{document} at the end of the file by removing the % from the start of each line.

%\documentclass[12pt, a4paper]{memoir}

%\input{LaTexPackages.tex}

%\begin{document}

%\begin{center}
	%\textbf{\large Your title goes here}

	%\textbf{Abstract}

	%Your Name, The University of Queensland, 20??
%\end{center}

% ********* Enter your text below this line: ********
Diffuse radio sources have been found within the cool cores of a growing number of relaxed galaxy clusters. These faint `minihalos' are a result of synchrotron-emitting, highly relativistic electrons. With extents over several hundred kiloparsec, they are far beyond the expected extent of a single source of relativistic electrons. The processes of \textit{in-situ} turbulent reacceleration and hadronic production have been proposed to explain this phenomenon. Spectral steepening is expected at high radio frequencies for the reacceleration scenario, however, the majority of previous studies have been performed at lower radio frequencies (144 MHz to 1.4 GHz) where this is difficult to detect.

To investigate potential steepening, in this thesis I calibrate and image Very Large Array radio interferometer observations at 10 GHz of 12 clusters hosting minihalos. This represents the first steps towards a high frequency survey of minihalos. Source-subtracted measurements of minihalo flux are performed for four clusters, enabling comparison with prior observations at lower frequencies. Steepening was not observed for three clusters, RXJ1720+2638 ($\alpha=1.1\pm0.05$), RXCJ1115+0129 ($\alpha=1.03\pm0.18$), and RXJ2129+0005 ($\alpha=1.18\pm0.09$). This suggests a hadronic origin for these minihalos. One cluster, A478 ($\alpha=2.02\pm0.10$), exhibits tentative evidence of steepening, however, more research is required to ensure this is not from instrumental effects. The remainder of the sample requires more extensive analysis. Challenges encountered during the calibration and imaging are discussed, and recommendations for future analysis of high frequency minihalo studies are presented. 
% ***************************************************

%\end{document}